\documentclass[11pt]{article}
\usepackage[margin=1.2in]{geometry}
\usepackage{amsmath,amssymb,amsthm}
\usepackage[hidelinks]{hyperref}


\newtheorem{theorem}{Theorem}
\newtheorem{proposition}[theorem]{Proposition}
\newtheorem{lemma}[theorem]{Lemma}
\newtheorem{corollary}[theorem]{Corollary}
\theoremstyle{definition}
\newtheorem{remark}[theorem]{Remark}

\newcommand{\SB}{\mathrm{SB}}
\newcommand{\pr}{\mathbb{P}}
\newcommand{\E}{\mathbb{E}}

\title{Ordered simplex caps and the three-observation Stringer bound}
\author{Richie Sater\\
\small Independent researcher, Silver Spring, Maryland, USA\\
\small \texttt{richiesater@gmail.com}}
\date{September 2026}

\begin{document}
\maketitle

\begin{abstract}
We prove that a cap of the uniform distribution on a tetrahedron, cut through an ordered
barycentric point, has maximal volume at a step normal within the matching
order cone. The proof reduces stationary caps to a section-centroid
inequality, whose only nonsimplicial case is established by exact rational
computation. Applied to the binomial Stringer bound for three i.i.d.\ observations in $[0,1]$, the theorem gives distribution-free coverage at
every confidence of at least $66.6314788\ldots\%$. This range is sharp for
pointwise domination of Gaffke's valid mean bound. In every sample size, a
two-knot perturbation gives an explicit obstruction to that comparison;
its critical tail levels converge to $0.2846681370\ldots$.
\end{abstract}

\noindent\textbf{Keywords:} Stringer bound; bounded mean; simplex cap;
Dirichlet average.

\section{Introduction}

The Stringer bound \cite{stringer1963} estimates the mean of a distribution
on $[0,1]$ from independent observations. In monetary-unit sampling these
observations are taints, the fractions by which sampled account items are
overstated. For a sample arranged as $t_{(1)}\ge\cdots\ge t_{(n)}$, the
binomial-factor bound at confidence $1-\alpha$, where $0<\alpha<1$, is
\begin{equation}\label{eq:sb}
 \SB=p_n(0)+\sum_{j=1}^n\{p_n(j)-p_n(j-1)\}t_{(j)},
\end{equation}
where, for $j<n$, the Clopper--Pearson factor $p_n(j)$ is the unique solution
in $(0,1)$ of
\[
 \sum_{i=0}^j\binom ni p^i(1-p)^{n-i}=\alpha,
 \qquad p_n(n)=1.
\]
Its finite-sample conservatism asks whether
\begin{equation}\label{eq:conj}
 \pr_F\{\SB\ge\mu(F)\}\ge1-\alpha
\end{equation}
for every distribution $F$ on $[0,1]$, with mean $\mu(F)$. The general
question at conventional confidence levels remains open. Bickel
\cite{bickel1992} established asymptotic conservatism; Pap and van Zuijlen
\cite{papvz1996} identified the asymptotic confidence threshold of $1/2$,
and their finite-sample examples \cite{papvz1995} show that low-confidence
failure is possible. Bernoulli taints reduce the bound to Clopper--Pearson
\cite{djpvz1997}. The present paper concerns independent draws, not
systematic or other dependent sampling designs.

Our main result is geometric. Write $I_z(a,b)$ for the distribution
function at $z$ of a $\operatorname{Beta}(a,b)$ random variable. A
$\operatorname{Dirichlet}(1,\ldots,1)$ vector is uniform on the simplex
$\{d_i\ge0:\sum_i d_i=1\}$.

\begin{theorem}[Four-coordinate monotone cap theorem]\label{thm:fourcap}
Let $D\sim\operatorname{Dirichlet}(1,1,1,1)$, let
$c_0\le c_1\le c_2\le c_3$ be nonnegative with $\sum_i c_i=1$, and let
$x_0\le x_1\le x_2\le x_3$.  Put $C_r=\sum_{i<r}c_i$.  Then
\begin{equation}\label{eq:fourcaptheorem}
 \Pr\{x\mathbin{\cdot}D>c\mathbin{\cdot}x\}
 \le\max_{r=1,2,3}I_{C_r}(r,4-r).
\end{equation}
Moreover, the maximum over all nonconstant ordered $x$ equals the right side:
each candidate is attained by the step vector with its first $r$ coordinates
zero and its remaining coordinates one.
\end{theorem}


Thus a cap cut through the barycentric point $c$ is largest when the
normal $x$ takes only two values. The assertion holds for every ordered
$c$, not only for confidence-bound coefficients.

For observations $t_1,\ldots,t_n\in[0,1]$, define
\begin{equation}\label{eq:gaffke}
 G_\alpha(t)=Q_{1-\alpha}\!\left(D_0+\sum_{i=1}^n t_iD_i\right),
 \qquad D\sim\operatorname{Dirichlet}(1,\ldots,1),
\end{equation}
where $Q_q(X)=\inf\{z:\pr(X\le z)\ge q\}$. Gaffke's mean test
\cite{gaffke2005}, proved valid by Vlassis and Thomas
\cite[Theorem~1]{vlassis2026}, gives
$\pr_F\{G_\alpha\ge\mu(F)\}\ge1-\alpha$.
The quantile formulation follows by test inversion
\cite[Theorem~5.1 and Corollary~5.2]{ming2026}.
The geometric theorem yields the following exact comparison range.

\begin{corollary}[Three-observation Stringer range]\label{cor:n3expanded}
Set
\begin{equation}\label{eq:alphastar}
 \alpha_\star=\left(\frac{19+\sqrt{21}}{34}\right)^3
 =0.3336852118672717\ldots .
\end{equation}
For $n=3$, the binomial Stringer bound dominates $G_\alpha$ on every
sample if and only if $0<\alpha\le\alpha_\star$. Consequently
\eqref{eq:conj} holds for every $F$ at every nominal confidence of at least
$66.6314788\ldots\%$.
\end{corollary}

The endpoint concerns pointwise comparison, not the coverage threshold of
Stringer itself. Proposition~\ref{prop:terminalbifurcation} gives the
corresponding obstruction in every dimension.

The obstacle is that a simplex cap moves both its normal and its threshold:
ordinary stochastic ordering of Dirichlet averages does not directly
compare these probabilities. The Dirichlet-average and spline formulas
are classical \cite{carlson1991,lasserre2015}. The centroid case is related
to the simplex halfspace-depth calculation of Paindaveine and Van Bever
\cite{paindaveine2017}; here the barycentric point is arbitrary and the
normal is restricted to its order cone. The correspondence
between stationary cuts and section centroids
\cite[Proposition~1.11]{patakova2022} is also classical. These facts calculate a cap and locate
its stationary points, but do not identify its global maximum in an order
cone.

The key additional inequality bounds a cap by a beta probability formed
from its section centroid. On a face where knots coincide, stationarity
identifies the centroid's block masses with those of $c$. It is therefore
essential to use prefixes ending at strict knot gaps. The resulting local
inequality bounds every face maximum by a step cap. For a tetrahedron,
the only nonelementary section has two knots on each side of the cut.

In that case, centroid order provides a polynomial constraint, and a
nonnegative Bernstein multiplier proves the required inequality. A separate
identity treats the repeated-upper-knot boundary. The finite calculation
uses exact rational arithmetic; the stationary-face reduction, the
Clopper--Pearson weight inequalities, and the limiting obstruction are
ordinary proofs.

\section{Section centroids and tetrahedral caps}\label{sec:fourcap}

For $N\ge2$, let $D\sim\operatorname{Dirichlet}(1,\ldots,1)$, let
$x_0\le\cdots\le x_{N-1}$, put $T=x\mathbin{\cdot}D$, and consider a
regular section $T=q$.  Write
\begin{equation}\label{eq:sectioncentroid}
 p=\Pr(T>q),\qquad
 \gamma_i=\E(D_i\mid T=q),\qquad
 M_r=\sum_{i<r}\gamma_i.
\end{equation}
At the step vector with its first $r$ knots zero and its remaining knots one,
the cap at the barycentric threshold $c\mathbin{\cdot}x$ has value
$I_{C_r}(r,N-r)$, where $C_r=\sum_{i<r}c_i$.

\begin{lemma}[Section-centroid maximum principle]\label{lem:sectionmaximum}
Suppose that, in dimension $N$, every regular ordered section with
$\gamma_0\le\cdots\le\gamma_{N-1}$ satisfies
\begin{equation}\label{eq:localbarrier}
 p\le \max_{r:\,x_{r-1}<x_r} I_{M_r}(r,N-r).
\end{equation}
Assume the same inequality on singular sections by continuity.  Then, for
every nondecreasing probability vector $c=(c_0,\ldots,c_{N-1})$, every
ordered $x$, and $C_r=\sum_{i<r}c_i$,
\begin{equation}\label{eq:globalcap}
 \Pr\{x\mathbin{\cdot}D>c\mathbin{\cdot}x\}
 \le \max_{1\le r<N} I_{C_r}(r,N-r).
\end{equation}
\end{lemma}

\begin{proof}
For constant $x$ the cap is empty. Otherwise, first take every $c_i>0$;
then $q=c\mathbin{\cdot}x$ lies strictly between the endpoint knots.
Translation and positive scaling of $x$ leave the cap unchanged, so normalize
its least and greatest distinct knots to zero and one.  On any face of this
compact ordered-knot polytope, collect the equal consecutive knots into
blocks.  If their distinct values are $y_1<\cdots<y_k$, put
$W_b=\sum_{i\text{ in block }b}D_i$ and let $A_b$ be the corresponding sum
of the $c_i$.  Differentiation of the cap at a regular section gives
\begin{equation}\label{eq:capderivative}
 \frac{\partial}{\partial y_b}
 \Pr\!\left\{\sum_a y_aW_a>\sum_a y_aA_a\right\}
 =f(q)\{\E(W_b\mid T=q)-A_b\},
\end{equation}
where $f$ is the density of $T$. At an interior knot level this identity
extends by continuity of the section integrals.

At a relative-interior stationary point, the right side vanishes for every
nonendpoint block.  The two identities
\[
 \sum_b\{\E(W_b\mid T=q)-A_b\}=0,
 \qquad
 \sum_b y_b\{\E(W_b\mid T=q)-A_b\}=0
\]
then give the same conclusion for the endpoint blocks.  Conditional symmetry
within a block makes each $\gamma_i$ its block mass divided by the block
size.  These values are consecutive block averages of the nondecreasing
vector $c$, and hence are nondecreasing.  Moreover, at every strict knot gap,
$M_r=C_r$.  The local barrier \eqref{eq:localbarrier} therefore bounds the
value at every relative-interior stationary point by the right side of
\eqref{eq:globalcap}.  A nonstationary face maximum lies on a proper face;
induction over the faces ends at a step vector, where the cap equals
$I_{C_r}(r,N-r)$.  Approximation and continuity cover the singular and
nonstrict cases.
\end{proof}

The first section with a nonsimplicial cross-section occurs when $N=4$ and
two knots lie on either side of the threshold.  The next proposition is the
local barrier needed in Lemma~\ref{lem:sectionmaximum}.

\begin{proposition}[Tetrahedral section barrier]\label{prop:tetrahedralbarrier}
Let $N=4$ and use the notation in \eqref{eq:sectioncentroid}.  If the section
centroid is ordered, then
\begin{equation}\label{eq:tetrahedralbarrier}
 p\le\max_{r:\,x_{r-1}<x_r}I_{M_r}(r,4-r).
\end{equation}
In the two-versus-two crossing $x_0<x_1<q<x_2<x_3$, the stronger inequality
\[
 p\le I_{M_2}(2,2)=3M_2^2-2M_2^3
\]
holds.
\end{proposition}

\begin{proof}
There are three crossing types.  If only $x_3$ is above $q$, set
$a_i=(x_3-q)/(x_3-x_i)$ for $i<3$.  The cap is a tetrahedron and its section
is its base, so
\[
 p=a_0a_1a_2,\qquad M_3=(a_0+a_1+a_2)/3.
\]
Thus $p\le M_3^3=I_{M_3}(3,1)$ by AM--GM.  If only $x_0$ is below $q$, set
$b_i=(q-x_0)/(x_i-x_0)$ for $i>0$.  Then
\[
 p=1-b_1b_2b_3,\qquad
 \gamma_0=1-(b_1+b_2+b_3)/3,\qquad \gamma_i=b_i/3.
\]
The knot order gives $b_1\ge b_2\ge b_3$, while the centroid order gives the
reverse inequalities.  Hence all three equal a common $b$, and
$p=1-b^3=I_{M_1}(1,3)$.

It remains to treat two knots on each side.  Translation and scaling put the
centered knots in the form
\begin{equation}\label{eq:normalizedfourknots}
 (-1,-u,tv,t),\qquad 0<u,v<1,\quad t>0.
\end{equation}
For positive temporary variables $A,B,C,D$, the divided-difference formula \cite{carlson1991,lasserre2015}
and its derivative give
\begin{align}
 p(A,B,C,D)&=
 \frac{C^3}{(C+A)(C+B)(C-D)}
 +\frac{D^3}{(D+A)(D+B)(D-C)},\label{eq:fourcapformula}\\
 \gamma&=f(0)^{-1}(-p_A,-p_B,p_C,p_D),
\end{align}
where the density $f(0)$ is obtained from \eqref{eq:fourcapformula} by
replacing $C^3,D^3$ with $3C^2,3D^2$.  This is followed by
$(A,B,C,D)=(1,u,tv,t)$.  Exact simplification of the upper
centroid-order condition gives
\begin{equation}\label{eq:orderfactor}
 \gamma_3-\gamma_2=\frac{(v-1)R(u,v,t)}{\Delta_0(u,v,t)},
 \qquad \Delta_0>0,
\end{equation}
where
\begin{align*}
 R={}&t^4v^2(u^2+u+1)+2t^3uv(u+1)(v+1)\\
    &+t^2u^2(v^2+4v+1)-u^3.
\end{align*}
Since $v<1$, orderedness gives $R\le0$.  The last positive term in $R$
also gives $t^2(v^2+4v+1)\le u<1$, so \eqref{eq:normalizedfourknots} lies
in the unit cube.

With $M=\gamma_0+\gamma_1$, exact rational simplification gives
\begin{equation}\label{eq:fourcomparison}
 3M^2-2M^3-p=\frac{t^2P(u,v,t)}{\Delta(u,v,t)},
\end{equation}
where $P$ has tensor degree $(9,9,11)$ and
\[
 \Delta=27(t+1)^3(t+u)^3(tv+1)^3(tv+u)^3
 (tuv+tv+uv+u)^3>0.
\]
The remaining inequality follows from an exact Bernstein-multiplier
identity. Recall that the tensor products of
$B_i^d(z)=\binom di z^i(1-z)^{d-i}$ are nonnegative on a unit cube;
thus nonnegative coefficients in this basis imply polynomial
nonnegativity.  There is a polynomial $S$ of tensor degree
$(6,7,7)$ with 79 positive Bernstein coefficients and all others zero.  Exact integer and rational calculation shows that
\[
 H=P+RS
\]
has 631 positive Bernstein coefficients of degree $(9,9,11)$ and all 569
remaining coefficients zero. The rational coefficients of $S$, and the
integer-arithmetic derivation and sign calculation, are supplied with the
code cited below. Hence on $R\le0$,
$P=H-RS\ge0$, proving \eqref{eq:fourcomparison}.

When $v=1$, equality of the upper knots removes the information in
\eqref{eq:orderfactor}.  The remaining centroid inequality has positive
denominator and numerator
\[
 G=-2t^3u-t^3+t^2u^2-4t^2u+3tu^2+3u^2.
\]
Orderedness gives $G\ge0$.  If $t\ge1$, then
\[
 G\le(u-1)\{(3u+1)t+4u\}\le0,
\]
so the nondegenerate ordered domain again has $t<1$.  After the positive
factor $t^2(1-u)^2$ is removed from the comparison numerator, the boundary
certificate checks the exact identity
\[
 P_\partial=H_\partial+2G(1-u)^2t^2,
\]
where all 30 degree-$(4,5)$ Bernstein coefficients of $H_\partial$ are
nonnegative (19 are positive, 11 are zero).  This closes the repeated-upper
boundary. The identity for $v<1$ remains valid at $u=1$, so
repeated lower knots cause no further case. If $u=v=1$, there are two
knot blocks and the cap equals $I_{M_2}(2,2)$. In the mixed crossing
$r=2$ is always a strict gap; in the endpoint crossings the active gap is
respectively $r=3$ or $r=1$. For an interior knot level $q=x_2<x_3$, the one-above formulas still
hold with $a_2=1$. If $x_0<q=x_1<x_2$, the section is a triangle with
$\gamma_1=1/3$ and $\gamma_2=(q-x_0)/(3(x_2-x_0))<1/3$,
contrary to centroid order. These exhaust the interior knot levels,
including $x_1=x_2$ in the first case.
\end{proof}

\begin{proof}[Proof of Theorem~\ref{thm:fourcap}]
Proposition~\ref{prop:tetrahedralbarrier} proves the hypothesis of
Lemma~\ref{lem:sectionmaximum} in every crossing stratum for $N=4$, yielding
\eqref{eq:fourcaptheorem}.  At the stated step vector,
$\sum_{i<r}D_i\sim\operatorname{Beta}(r,4-r)$ and the cap event is
$\sum_{i<r}D_i<C_r$, which proves attainment.
\end{proof}

\section{The sharp Stringer comparison}

Affine normalization turns comparison with \eqref{eq:gaffke} into the
cap problem above. The normalization is useful for any sample size.

\begin{proposition}[Simplex-cap criterion]\label{prop:capcriterion}
Arrange a sample as $t_1\le\cdots\le t_n$. If $t_1<1$, put
\[
z_0=0,\qquad z_k=\frac{t_{k+1}-t_1}{1-t_1}\ (1\le k\le n-1),
\qquad z_n=1,
\]
and
\[
s_{n,\alpha}(z)=p_n(0)+\sum_{k=1}^{n-1}
 \bigl(p_n(n-k)-p_n(n-k-1)\bigr)z_k.
\]
Then $\SB\ge G_\alpha(t_1,\ldots,t_n)$ if and only if
\begin{equation}\label{eq:capcriterion}
 \Pr_D\!\left\{\sum_{k=1}^{n-1}z_kD_k+D_n
        >s_{n,\alpha}(z)\right\}\le\alpha,
 \qquad 0\le z_1\le\cdots\le z_{n-1}\le1.
\end{equation}
Consequently, proving \eqref{eq:capcriterion} on the complete ordered
knot simplex proves distribution-free Stringer coverage at $(n,\alpha)$.
\end{proposition}

\begin{proof}
Rewrite \eqref{eq:sb} for ascending observations and subtract $t_1$ from
both the Stringer threshold and the $n+1$ knots
$(t_1,\ldots,t_n,1)$ in the Dirichlet average. Division by $1-t_1$
gives the displayed threshold and knots. Exchangeability of the Dirichlet
weights permits the stated labeling. Because the normalized average has
knots zero and one, its distribution is continuous, so its
$(1-\alpha)$-quantile is at most the threshold exactly when its strict
upper tail there is at most $\alpha$. If $t_1=1$, every observation and
both bounds equal one. The coverage conclusion follows from the validity
of $G_\alpha$.
\end{proof}

\begin{proof}[Proof of Corollary~\ref{cor:n3expanded}]
Write $x=p_3(0)$, $y=p_3(1)$, and $z=p_3(2)$.  In the ascending-knot form of
Proposition~\ref{prop:capcriterion}, Stringer's threshold has weight vector
\[
 c=(1-z,z-y,y-x,x).
\]
At the three step vectors, the beta--binomial identity makes the corresponding
cap probability exactly $\alpha$.  Theorem~\ref{thm:fourcap} therefore proves
the cap criterion once $c$ is nondecreasing.

Put
\[
 f(s)=1-3s^2+2s^3,\qquad f(y)=\alpha,\qquad
 x=1-\alpha^{1/3},\qquad z=(1-\alpha)^{1/3}.
\]
The function $f$ is strictly decreasing on $(0,1)$.  The last weight
inequality is $y\le2x$.  It is automatic if $2x\ge1$; otherwise it is
equivalent to
\begin{equation}\label{eq:lastweight}
 \alpha-f(2x)=x(-3+15x-17x^2)\ge0.
\end{equation}
The first root of the quadratic factor is
$x_\star=(15-\sqrt{21})/34$, and $\alpha=(1-x)^3$, so
the last weight inequality holds exactly through $\alpha_\star$.

For the middle inequality, put $r=\alpha^{1/3}=1-x$.  Since
$\alpha\le\alpha_\star<1/2$, we have $z\ge r$ and $z^3+r^3=1$, while
\[
 4\left\{f\left(\frac{x+z}{2}\right)-\alpha\right\}
 =3(z-r)(z^2+r^2-1)\ge0.
\]
Thus $y\ge(x+z)/2$.  Finally, $z>3/4$ and
\[
 \alpha-f(2z-1)=(1-z)(17z^2-19z+5)>0,
\]
so $y\le2z-1$.  These are precisely the three inequalities saying that $c$
is nondecreasing.  The cap criterion and finite-sample validity of $G_\alpha$
prove the first assertion.

For sharpness, suppose $\alpha>\alpha_\star$.  Then
$x<x_\star<1/2$, and \eqref{eq:lastweight} has the opposite strict sign;
since $f$ is decreasing, $y>2x$.  Thus the last two threshold weights satisfy
$c_2=y-x>c_3=x$.  Consider the ordered knot family
$k(h)=(0,0,h,1)$ and write
\[
 \Phi(h)=\Pr\{k(h)\mathbin{\cdot}D>c\mathbin{\cdot}k(h)\}.
\]
At $h=1$, the beta--binomial identity gives $\Phi(1)=\alpha$.
Conditional on $D_2+D_3=c_2+c_3$, symmetry gives
$\E(D_2\mid D_2+D_3=c_2+c_3)=(c_2+c_3)/2$.  The derivative formula
\eqref{eq:capderivative} therefore yields
\[
 \Phi'(1)=\frac{g_\alpha}{2}(c_3-c_2)<0,
\]
where $g_\alpha>0$ is the density of $D_2+D_3$ at $c_2+c_3$.
Consequently $\Phi(h)>\alpha$ for $h<1$ sufficiently close to one, and
Proposition~\ref{prop:capcriterion} gives $\SB<G_\alpha$ on the corresponding
sample.
\end{proof}

\section{A terminal-edge obstruction in every dimension}

The same perturbation identifies a necessary condition for the comparison
at every sample size.

\begin{proposition}[Terminal-edge phase transition]
\label{prop:terminalbifurcation}
For each $n\ge2$, let $r_n$ be the unique solution in
$(1/2,1-1/(2n))$ of
\begin{equation}\label{eq:terminalcritical}
 (2r-1)^{n-1}\{2n-1-2(n-1)r\}=r^n,
\end{equation}
and put $\alpha_n=r_n^n$.  If $\alpha>\alpha_n$, then some ordered sample
satisfies $\SB<G_\alpha$.  Moreover,
\begin{equation}\label{eq:terminalcritical-limit}
 \alpha_n\longrightarrow\alpha_\infty
 =0.28466813704083846168\ldots,
 \qquad \alpha_\infty(1-2\log\alpha_\infty)=1.
\end{equation}
Thus for every $\alpha>\alpha_\infty$, pointwise Stringer--Gaffke
domination fails for all sufficiently large $n$.
\end{proposition}

\begin{proof}
Write $p_0=p_n(0)$ and $p_1=p_n(1)$, and consider
$x(\varepsilon)=(0^{\,n-1},1-\varepsilon,1)$.  If
$D\sim\operatorname{Dirichlet}(1,\ldots,1)$, then
\[
 S=D_{n-1}+D_n\sim\operatorname{Beta}(2,n-1),\qquad
 W=D_{n-1}/S\sim\operatorname{Unif}(0,1)
\]
independently.  The random score is $S(1-\varepsilon W)$ and its barycentric
threshold is $p_1-\varepsilon(p_1-p_0)$.  Differentiating the conditional
beta survival probability at zero gives
\begin{equation}\label{eq:terminalderivative}
 \left.\frac{d}{d\varepsilon}\Phi_c(x(\varepsilon))
 \right|_{\varepsilon=0}
 =f_{2,n-1}(p_1)\left(\frac{p_1}{2}-p_0\right).
\end{equation}
Here $\Phi_c(x)=\pr\{x\mathbin{\cdot}D>c\mathbin{\cdot}x\}$,
and $f_{2,n-1}$ is the beta density. The cap at zero is the calibrated
step cap of probability $\alpha$.
Consequently it rises above $\alpha$ exactly when $p_1>2p_0$; equivalently,
the last two threshold weights reverse order because
$c_{n-1}=p_1-p_0>p_0=c_n$.  The cap criterion then gives
$\SB<G_\alpha$ for the corresponding sample.

Put $r=\alpha^{1/n}$, so $p_0=1-r$.  Substituting $p_1=2p_0$ in its
defining binomial tail gives \eqref{eq:terminalcritical}.  For
\[
 H_n(r)=\frac{(2r-1)^{n-1}\{2n-1-2(n-1)r\}}{r^n},
\]
direct differentiation yields
\[
 \frac{d}{dr}\log H_n(r)
 =\frac{n(2nr-2n+1)}
 {r(2r-1)(2nr-2n-2r+1)}.
\]
The last denominator factor is negative.  Hence $H_n$ increases to
$r=1-1/(2n)$ and then decreases strictly to $H_n(1)=1$, proving the unique
nontrivial crossing and the stated direction of its sign.  Finally, for
fixed $\alpha$,
\[
 H_n(\alpha^{1/n})\longrightarrow
 \alpha(1-2\log\alpha).
\]
The limiting function has exactly one nontrivial root in $(0,1)$; the
convergence is locally uniform around that root, so the unique finite roots
converge to it.  This proves \eqref{eq:terminalcritical-limit}
and the eventual-failure assertion.
\end{proof}


For $n=3$, this critical level is exactly $\alpha_\star$.
The derivative identifies why comparison fails: the last two barycentric
weights reverse order. It does not construct an i.i.d. distribution with
Stringer undercoverage. The tetrahedral theorem therefore suggests two
distinct questions: extending the active-gap section inequality to higher
dimensions, and proving coverage beyond the range of pointwise comparison.

\section*{Data and code availability}
The exact rational multiplier, the boundary identity, and their derivation
and sign checks are available in the accompanying
\href{https://github.com/RichieSater/stringer-bound}{source repository}.

\section*{AI disclosure}
Anthropic Claude and OpenAI Codex were used for mathematical exploration,
code development, literature searches, and drafting; the author accepts
responsibility for the manuscript.

\noindent The author received no external funding and declares no competing
interests.

\begin{thebibliography}{99}

\bibitem{carlson1991} B.~C. Carlson, B-splines, hypergeometric functions,
and Dirichlet averages, \emph{Journal of Approximation Theory} \textbf{67}
(1991), 311--325,
\url{https://doi.org/10.1016/0021-9045(91)90006-V}.

\bibitem{bickel1992} P.~J. Bickel, Inference and auditing: the Stringer
bound, \emph{International Statistical Review} \textbf{60} (1992),
197--209.

\bibitem{djpvz1997} N.~G. de~Jager, G.~Pap and M.~C.~A. van~Zuijlen,
Facts, phantasies, and a new proposal concerning the Stringer bound,
\emph{Computers \& Mathematics with Applications} \textbf{33} (1997),
37--54.

\bibitem{gaffke2005} N.~Gaffke, Three test statistics for a nonparametric
one-sided hypothesis on the mean of a nonnegative variable,
\emph{Mathematical Methods of Statistics} \textbf{14} (2005), 451--467.

\bibitem{lasserre2015} J.~B. Lasserre, Volume of slices and sections of the
simplex in closed form, \emph{Optimization Letters} \textbf{9} (2015),
1263--1269, \url{https://doi.org/10.1007/s11590-015-0898-z}.

\bibitem{ming2026} J.~Ming, A.~Ramdas, Y.~Shen, R.~Wang and
I.~Waudby-Smith, Gaffke's confidence interval for the mean of bounded data
is inadmissible but asymptotically efficient, arXiv:2607.18661 (2026).

\bibitem{patakova2022} Z.~Pat\'akov\'a, M.~Tancer and U.~Wagner,
Barycentric cuts through a convex body, \emph{Discrete \& Computational
Geometry} \textbf{68} (2022), 1133--1154,
\url{https://doi.org/10.1007/s00454-021-00364-7}.

\bibitem{paindaveine2017} D.~Paindaveine and G.~Van Bever,
On the maximal halfspace depth of permutation-invariant distributions on
the simplex, \emph{Statistics \& Probability Letters} \textbf{129}
(2017), 335--339, \url{https://doi.org/10.1016/j.spl.2017.06.019}.

\bibitem{papvz1995} G.~Pap and M.~C.~A. van~Zuijlen, The Stringer bound
in case of uniform taintings, \emph{Computers \& Mathematics with
Applications} \textbf{29} (1995), 51--59.

\bibitem{papvz1996} G.~Pap and M.~C.~A. van~Zuijlen, On the asymptotic
behaviour of the Stringer bound, \emph{Statistica Neerlandica}
\textbf{50} (1996), 367--389.

\bibitem{vlassis2026} N.~Vlassis and P.~S. Thomas, An exact
distribution-free test for means of nonnegative random variables,
arXiv:2607.08415 (2026).

\bibitem{stringer1963} K.~W. Stringer, Practical aspects of statistical
sampling in auditing, \emph{Proceedings of the Business and Economic
Statistics Section, American Statistical Association} (1963), 405--411.

\end{thebibliography}
\end{document}
